\chapter{Introduction générale}

\section{Contexte général}

Les journaux système et réseau occupent une place centrale dans la supervision, l'audit et la cybersécurité. Ils enregistrent une grande diversité d'événements : connexions, erreurs, changements d'état, accès aux services, alertes de sécurité ou encore activités applicatives. Dans une infrastructure moderne, ces informations peuvent provenir de systèmes Windows, de serveurs Linux, d'applications web, d'équipements réseau, de plateformes SIEM, d'environnements cloud ou encore de jeux de données publics utilisés pour l'évaluation scientifique.

Cette diversité constitue toutefois une difficulté importante. Les journaux sont généralement volumineux, hétérogènes et parfois incomplets. Ils ne suivent pas non plus une structure commune. Un événement Windows, une ligne syslog Linux, une requête Apache, une alerte Wazuh ou un flux réseau représenté sous forme tabulaire ne décrivent pas nécessairement les mêmes informations, avec la même granularité temporelle ni le même niveau de détail. Dans ces conditions, l'analyse manuelle devient rapidement coûteuse et difficile à maintenir, surtout dans les petites structures, les établissements universitaires ou les organisations qui ne disposent pas d'un centre opérationnel de sécurité complet.

Les travaux consacrés à la détection d'intrusions ont depuis longtemps montré l'intérêt des traces d'audit et de l'identification des comportements inhabituels \cite{denning1987intrusion,axelsson2000intrusion}. Les approches plus récentes fondées sur la détection d'anomalies et l'apprentissage automatique offrent des possibilités supplémentaires. Elles restent cependant sensibles à plusieurs facteurs, notamment la qualité des données, le déséquilibre des classes, les changements de distribution et la présence de faux positifs \cite{chandola2009anomaly,buczak2016survey,khraisat2019survey}. Il devient dès lors pertinent de disposer d'une chaîne de traitement capable de collecter des journaux hétérogènes, de les normaliser, de les orienter vers des traitements adaptés et de présenter les résultats sous une forme exploitable.

\section{Problématique}

La question centrale étudiée dans ce mémoire est la suivante :

\begin{quote}
Comment concevoir un prototype autonome, distribué au sens logiciel, modulaire et exploitable par un analyste, capable de détecter et de prioriser des anomalies candidates dans des journaux système et réseau hétérogènes au moyen de traitements périodiques, de micro-lots et d'agents multitâches ?
\end{quote}

Cette problématique rassemble plusieurs difficultés. La première concerne la reconnaissance et la structuration de formats de journaux très différents. La deuxième concerne le choix du traitement ou du modèle le plus approprié à chaque famille de données. Il faut également transformer des scores techniques en informations suffisamment claires pour être interprétées par un analyste. Enfin, l'évaluation doit rester rigoureuse et éviter une confusion importante : une anomalie produite par un modèle n'est pas nécessairement une attaque confirmée.

\section{Définition opérationnelle du titre}

Le titre du mémoire est maintenu dans sa formulation retenue, mais les notions qu'il contient doivent être comprises dans un sens opérationnel et vérifiable.

Dans ce travail, la détection autonome signifie que le prototype est capable de sélectionner certains traitements, d'orienter les différentes familles de journaux et d'exécuter des tâches sans intervention manuelle constante. Cette autonomie repose sur des règles explicites et sur un registre de modèles. Elle ne doit donc pas être comprise comme une autonomie décisionnelle complète en matière de cybersécurité, encore moins comme une forme d'intelligence générale capable de remplacer l'analyste.

La notion de distribution renvoie à la décomposition des traitements en plusieurs agents logiciels. Dans la version avancée du prototype, une répartition locale multiprocessus est réalisée avec Redis Streams. Plusieurs workers peuvent consommer une file commune, publier leur progression et reprendre certaines tâches qui n'ont pas été acquittées. Cette expérimentation ne constitue pas encore une démonstration de scalabilité industrielle sur plusieurs machines. Elle montre néanmoins que l'architecture peut dépasser un fonctionnement strictement monolithique.

Les agents intelligents multitâches correspondent à des agents logiciels disposant de capacités déclarées, d'une mémoire locale, d'un mécanisme de heartbeat et d'une politique de sélection des tâches. La mémoire conserve notamment les succès, les erreurs, certaines durées d'exécution et les décisions récentes. Plusieurs handlers sont également intégrés pour assurer la découverte des sources, le parsing, le routage des modèles, la détection et la corrélation. Le terme \textit{intelligent} renvoie donc ici à une logique d'orchestration adaptative et d'ingénierie logicielle, et non à une cognition comparable à celle d'un être humain.

\section{Questions de recherche}

La problématique générale est décomposée en plusieurs questions de recherche :

\begin{enumerate}

\item Comment construire un schéma de normalisation capable de rapprocher des journaux système, réseau et SIEM sans perdre les informations nécessaires à l'analyse ?

\item Comment organiser plusieurs agents logiciels spécialisés de manière à conserver une architecture modulaire, auditable et extensible ?

\item Comment sélectionner automatiquement un modèle de détection adapté à la famille de données traitée ?

\item Comment présenter les anomalies candidates dans un tableau de bord permettant à l'analyste de les comprendre, de les filtrer et de les valider ?

\item Comment exploiter une mémoire persistante des décisions, des erreurs et des validations passées afin d'améliorer progressivement la priorisation et de préparer la réduction des faux positifs ?

\item Comment évaluer le prototype à partir de métriques quantitatives, de captures fonctionnelles, de mesures de latence et d'indicateurs de consommation des ressources ?

\end{enumerate}

\section{Hypothèses de travail}

Ce mémoire repose sur cinq hypothèses principales.

La première hypothèse est qu'il est possible de rapprocher des journaux hétérogènes au moyen d'un noyau commun de champs sans supprimer les informations originales utiles à l'audit.

La deuxième hypothèse concerne l'architecture. L'utilisation d'agents logiciels spécialisés pourrait améliorer la modularité, la traçabilité et la robustesse de l'exécution par rapport à une chaîne monolithique. Cette hypothèse est partiellement étudiée à travers la progression fonctionnelle du prototype, la campagne Redis locale de six heures et une comparaison contrôlée entre l'architecture monolithique et l'architecture à base d'agents. Cette comparaison mesure surtout le surcoût introduit par l'orchestration ainsi que la capacité de reprise après une panne. La modularité et la traçabilité reposent, quant à elles, principalement sur les choix de conception logicielle et sur les mécanismes d'audit mis en place.

La troisième hypothèse concerne le routage par famille de journaux. L'idée est qu'une sélection du traitement selon la nature de la source peut limiter les incompatibilités entre les données, les caractéristiques extraites et le modèle appliqué. Une première comparaison contrôlée entre un modèle global et des modèles spécialisés par famille est réalisée dans ce mémoire. Les performances restent cependant proches lorsque les variantes utilisent le même espace de caractéristiques. Les résultats obtenus ne permettent donc pas d'affirmer que le routage spécialisé est systématiquement supérieur.

La quatrième hypothèse est que des modèles relativement légers peuvent conserver un intérêt dans une chaîne locale et reproductible lorsqu'ils sont évalués au moyen de métriques clairement définies, associés à un mécanisme de corrélation et complétés par une validation humaine.

Enfin, la cinquième hypothèse suppose qu'une mémoire persistante des exécutions, des erreurs, des décisions et des validations de l'analyste peut soutenir une adaptation progressive du comportement des agents. Elle peut notamment contribuer à limiter certains traitements répétitifs, à mieux prioriser certaines sources et à préparer une réduction contrôlée des faux positifs. Cette mémoire n'est cependant pas assimilée à un mécanisme de réentraînement automatique et permanent des modèles. Elle correspond plutôt à une capitalisation vérifiable de l'historique opérationnel et des retours humains.

\section{Objectif général}

L'objectif général de ce travail est de concevoir et d'évaluer un prototype modulaire, appelé \logminer{}, destiné à détecter et à prioriser des anomalies candidates dans des journaux système et réseau hétérogènes.

Le prototype doit couvrir l'ensemble de la chaîne de traitement, depuis la collecte des journaux jusqu'à la visualisation des résultats, en intégrant la normalisation, le routage vers les modèles appropriés, la détection et la corrélation.

\section{Objectifs spécifiques}

Les objectifs spécifiques retenus conformément à notre cahier des charges sont les suivants :

\begin{enumerate}

\item Identifier, catégoriser et structurer les différents types de journaux système et réseau utiles à l'analyse de sécurité.

\item Étudier les principales techniques classiques et modernes de détection d'anomalies.

\item Concevoir une architecture multiagent extensible et expérimenter localement la distribution des tâches au moyen de Redis Streams.

\item Intégrer des mécanismes d'apprentissage destinés à la détection, notamment des modèles supervisés et non supervisés, ainsi qu'une mémoire persistante permettant d'intégrer certains retours de l'analyste.

\item Concevoir un tableau de bord destiné à la visualisation des alertes et des anomalies candidates.

\item Tester le prototype sur des journaux simulés et sur des données réelles ou issues de jeux de données publics.

\item Évaluer ses performances à l'aide de métriques telles que la précision, le rappel, le score F1, la latence et la consommation des ressources.

\end{enumerate}

\section{Contributions}

La contribution principale de ce mémoire est essentiellement architecturale. Le travail ne propose pas un nouvel algorithme de détection, une nouvelle fonction objectif ou une nouvelle théorie d'apprentissage. Il porte plutôt sur la conception et l'évaluation d'une architecture d'intégration capable d'orchestrer plusieurs traitements spécialisés sur des journaux hétérogènes.

Cette orientation est cohérente avec l'objectif général du travail : construire une chaîne modulaire, reproductible et techniquement vérifiable.

Une comparaison contrôlée entre un modèle global et un routage par famille est notamment menée sur cinq graines aléatoires dans un espace de caractéristiques commun. Cette comparaison est complétée par l'étude de plusieurs configurations spécialisées. Des expérimentations supplémentaires utilisant des partitions temporelles et plusieurs environnements resteraient néanmoins nécessaires pour approfondir cette analyse.

Les principales contributions du mémoire sont les suivantes :

\begin{itemize}

\item une architecture logicielle modulaire destinée à l'analyse de journaux hétérogènes ;

\item un schéma commun de normalisation permettant de rapprocher plusieurs sources de journaux ;

\item un mécanisme de routage multimodèle par famille de journaux afin d'éviter l'application systématique d'un modèle unique à toutes les sources ;

\item un environnement d'exécution d'agents multitâches accompagné d'une preuve locale de répartition des tâches à l'aide de Redis Streams ;

\item une mémoire adaptative auditable reliant l'historique des exécutions, les décisions de l'analyste et différents signaux de retour afin de préparer une amélioration progressive du comportement des agents ;

\item une évaluation réalisée sur plusieurs jeux de données, comprenant des journaux locaux, des exports Wazuh, des jeux de données publics et différents scénarios de robustesse ;

\item un tableau de bord web permettant de visualiser les résultats, de filtrer les événements, de valider certaines décisions et de conserver une trace de l'analyse ;

\item un ensemble d'artefacts reproductibles comprenant des figures, des tableaux, des captures et des fichiers issus des expérimentations.

\end{itemize}

\section{Délimitation du travail}

Le travail porte principalement sur la conception, l'implémentation et l'évaluation d'un prototype local reproductible.

La dimension distribuée est étudiée à travers une architecture multiagent et une expérimentation multiprocessus locale utilisant Redis Streams. Elle ne doit donc pas être interprétée comme la démonstration d'un déploiement industriel multisite complet.

De la même manière, les résultats issus des modèles supervisés décrivent les performances obtenues sur les jeux de données et les protocoles retenus. Ils ne démontrent pas que les mêmes performances seront automatiquement obtenues sur toutes les infrastructures réelles.

La mémoire adaptative est considérée comme un mécanisme d'amélioration comportementale et d'aide à la décision. Elle ne constitue pas encore un système d'apprentissage continu modifiant automatiquement les poids des modèles utilisés en production.

\section{Méthodologie générale}

La démarche suivie est à la fois itérative et expérimentale.

Elle commence par l'étude des formats de journaux et des principales méthodes utilisées pour détecter des anomalies. Cette analyse conduit ensuite à la conception d'une architecture modulaire, puis à l'implémentation progressive du prototype \logminer{}.

Les modèles intégrés sont ensuite évalués sur plusieurs familles de données. Les résultats produits sont regroupés dans des tableaux, des figures et des captures du tableau de bord.

Le prototype privilégie des modèles relativement légers pouvant fonctionner dans un environnement local. Les approches fondées sur l'apprentissage profond sont étudiées dans l'état de l'art, mais l'implémentation se concentre principalement sur des méthodes plus simples à intégrer et à reproduire, parmi lesquelles Isolation Forest, Random Forest, le routage par famille et les mécanismes de corrélation des anomalies candidates.

La validation ne repose pas sur une seule métrique. Elle combine des mesures numériques, des tableaux synthétiques, des figures exportées, des captures du tableau de bord, des scripts permettant de reproduire certaines expériences et des comparaisons encadrées avec des outils ou méthodes de référence. Cette approche permet d'évaluer non seulement les performances des modèles, mais également l'exploitabilité et la reproductibilité de l'ensemble du prototype.

\section{Organisation du mémoire}

Le mémoire reprend la structure définie dans le document directeur tout en ajoutant plusieurs parties consacrées à l'analyse scientifique, à la discussion et à la reproductibilité.

Le chapitre 2 présente l'état de l'art consacré aux journaux, aux méthodes de détection d'anomalies et aux architectures utilisées dans le domaine.

Le chapitre 3 décrit le modèle proposé ainsi que l'architecture modulaire de \logminer{}.

Le chapitre 4 présente l'implémentation du prototype et les principaux composants logiciels.

Le chapitre 5 est consacré aux résultats expérimentaux et à leur analyse.

Le chapitre 6 approfondit certains aspects scientifiques du travail, notamment les menaces à la validité et plusieurs choix méthodologiques.

Le chapitre 7 discute les résultats obtenus, les limites du prototype et sa valeur pratique.

Le chapitre 8 présente les éléments de reproductibilité, le dépôt GitHub, les scénarios de déploiement et certaines précautions d'exploitation.

Le dernier chapitre reprend la conclusion générale ainsi que les principales perspectives.

Enfin, les annexes regroupent les tableaux, les captures, les figures et plusieurs éléments permettant de vérifier ou de reproduire les résultats présentés dans le mémoire.

\section{Motivation du choix du sujet}

Le choix de ce sujet repose sur un constat relativement simple : presque toutes les infrastructures informatiques produisent des journaux, mais ces données ne sont pas toujours exploitées de manière systématique.

Une machine Windows génère par exemple des événements liés à la sécurité, au système et aux applications. Un serveur Linux conserve différentes traces d'authentification et d'activité des services. Un serveur web enregistre les requêtes reçues et les erreurs rencontrées. Un outil HIDS ou SIEM comme Wazuh ajoute quant à lui ses propres alertes, règles et métadonnées.

Malgré cette quantité d'informations, les différentes sources restent souvent séparées. Elles peuvent être consultées uniquement lorsqu'un incident se produit ou être analysées indépendamment les unes des autres.

Cette situation présente à la fois un intérêt scientifique et un intérêt pratique. Les journaux peuvent contenir des signaux faibles révélateurs d'un comportement inhabituel, mais leur volume et leur hétérogénéité rendent leur exploitation manuelle difficile.

Les outils existants fournissent déjà une aide importante, notamment grâce aux règles et aux signatures. Les méthodes d'apprentissage automatique peuvent compléter ces mécanismes en identifiant certains comportements atypiques. Leur utilisation reste toutefois pertinente uniquement si elles sont intégrées dans une chaîne suffisamment claire, contrôlable et reproductible.

Ce mémoire se situe donc entre deux extrêmes : d'un côté, le simple filtrage manuel de journaux ; de l'autre, une plateforme SOC industrielle complexe. Le prototype vise à proposer une solution plus structurée et plus analytique que la première, tout en restant plus légère et plus facilement vérifiable que la seconde.

Cette motivation est également liée au contexte local. Une université ou une petite organisation ne dispose pas nécessairement des moyens nécessaires pour déployer immédiatement une infrastructure complète de supervision et de réponse aux incidents.

Il reste cependant possible de construire progressivement une solution capable de collecter les journaux disponibles, de les structurer, de les analyser et de présenter les résultats. Cette démarche permet d'identifier les besoins réels, de mieux comprendre les différentes sources, de mesurer les limites des traitements utilisés et de préparer une éventuelle évolution vers une infrastructure plus complète.

\section{Intérêt scientifique et intérêt d'ingénierie}

L'intérêt scientifique de ce travail repose sur plusieurs dimensions complémentaires.

La première concerne la normalisation de formats différents. Sans représentation commune, il devient difficile de comparer les sources et de les orienter correctement vers les traitements appropriés.

La deuxième concerne le routage par famille. Le prototype ne considère pas qu'un modèle unique soit nécessairement adapté à tous les types de journaux. Il cherche plutôt à sélectionner un traitement en fonction de la nature de la source.

La troisième dimension concerne l'utilisation conjointe de méthodes supervisées et non supervisées. Lorsque les données sont labellisées, les performances peuvent être mesurées à l'aide de métriques classiques. Lorsque les labels ne sont pas disponibles, les modèles non supervisés permettent plutôt d'identifier des anomalies candidates qui doivent ensuite être interprétées avec prudence.

La quatrième dimension porte sur l'intégration de mécanismes de visualisation et d'audit. Le tableau de bord permet de rapprocher les sorties techniques du système de la décision humaine et de conserver une trace des actions réalisées.

L'intérêt d'ingénierie réside surtout dans la réalisation concrète du prototype. Le travail ne se limite pas à une revue de littérature ou à quelques expérimentations réalisées dans un notebook. Il aboutit à un dépôt organisé comprenant des modules Python, des scripts de collecte et de benchmark, des modèles sauvegardés, des figures, des tableaux et une interface web.

Cette matérialisation est importante dans le cadre d'un mémoire d'ingénieur civil. Elle montre le passage d'un problème identifié à une architecture, puis de cette architecture à une implémentation pouvant être exécutée, observée et mesurée.

Le mémoire cherche néanmoins à conserver un positionnement raisonnable. \logminer{} n'est présenté ni comme un produit industriel terminé ni comme un détecteur universel. Il s'agit d'une chaîne prototypée, testée et documentée, capable de traiter plusieurs familles de journaux et de fournir des résultats pouvant être étudiés dans un cadre scientifique.

\section{Contraintes et choix de conception}

Plusieurs contraintes ont influencé les choix réalisés pendant le développement.

La première concerne la légèreté du système. Le prototype devait pouvoir fonctionner localement sans nécessiter une infrastructure particulièrement lourde. Cette contrainte a notamment orienté le choix vers Python, scikit-learn, FastAPI, les formats CSV et JSONL ainsi qu'une interface web.

La deuxième contrainte concerne la reproductibilité. Les valeurs présentées dans le mémoire devaient pouvoir être reliées à des scripts, à des fichiers de résultats ou à d'autres artefacts techniques. L'objectif était d'éviter une évaluation uniquement déclarative.

Une troisième difficulté concerne la disponibilité des données. Les journaux issus d'environnements réels ne sont pas toujours complets, accessibles, partageables ou labellisés. Le projet combine par conséquent plusieurs sources : journaux locaux, exports Wazuh, jeux de données publics et scénarios contrôlés. Cette combinaison ne remplace pas une évaluation menée directement dans un SOC industriel, mais elle permet de disposer d'une base expérimentale plus diversifiée.

Enfin, une contrainte importante concerne la rigueur scientifique dans l'interprétation des résultats. Les sorties des modèles non supervisés sont donc présentées comme des anomalies candidates. Un déploiement distribué sécurisé sur plusieurs machines reste considéré comme une perspective, tout comme un véritable apprentissage continu intégrant un réentraînement contrôlé des modèles.

La mémoire adaptative est, quant à elle, limitée à un mécanisme de capitalisation de l'historique des exécutions et des retours disponibles dans le prototype.

Ces contraintes expliquent plusieurs décisions prises dans le reste du mémoire : le choix de modèles relativement légers, l'intégration du tableau de bord dans la contribution, la conservation de fichiers de preuve dans les annexes et la discussion explicite des limites expérimentales.
